% Section 2 of the CIRP CMS paper: related work. Draft of 3 Nov 2026.
%
% Written from docs/paper-outline.md (section 2) and from the full texts of the works marked "full
% text" below (in Papers-cirp-cms, not in the repository; read 3 Nov). For rooker2007zero we have the
% abstract only. Budget: about 850 words and Table 1; with section 1 it fills the first two pages. The
% paper cites at most 17 sources: those of references.bib beside this file. Left out to keep to 17 (in
% docs/paper-references.bib): sidorenko2024composition, wenger2018connecting, lepuschitz2011self,
% koecher2024smt, bernhard2024leverage, nabizada2024workflow, css2022, strasser2008framework,
% jensen2026ontology. Needs amssymb (\checkmark), which the template loads.

\section{Related work}
\label{sec:related}

Four lines of work meet in this paper: composing skills on controllers, the AAS as a source for
control, changing IEC~61499 programs online, and planning from capabilities. Table~\ref{tab:related}
compares the closest works.

% Full text: dorofeev2021orchestrator, terzimehic2021service, lindorfer2019user.
\emph{Composing skills on controllers.} A skill encapsulates a device's function behind a uniform
interface, and a station's skills are composed from its devices'
primitives~\cite{schou2018skill}; the capability a skill realizes states what it achieves,
independent of how~\cite{koecher2023reference}. Dorofeev et al. generate a composite skill as
IEC~61499 code from a given orchestration sequence and from skill interfaces modelled as
adapters~\cite{dorofeev2021orchestrator}. This cut the engineering effort in 4diac by a factor of 2.7
to 5.5, but the result is one basic function block whose state machine grows to $2^n$ states for $n$
skills called in parallel. The skill descriptions are part of the IEC~61499 application; keeping them
in the AAS is the authors' future work. Terzimehić et al. compute where services are deployed with
SMT, from a recipe, the plant topology and the component models, derive the orchestration sequence
and generate an IEC~61499 application from it~\cite{terzimehic2021service}. Lindorfer and Froschauer
let end users model skill sequences in BPMN and validate each model against a rule set, so that an
invalid sequence cannot run~\cite{lindorfer2019user}. The sequence is executed by a workflow engine
that calls the skills over OPC~UA; running it on 4diac FORTE is their future work. In all of these, a
composition is made at engineering time and either deployed as new code or interpreted outside the
controller's program. None reads the deployed program back to compare it with the description it
came from.

% Full text: blumhofer2025enhancing, stolze2024proactive, albus2025automated. Cut first: "Our earlier
% work ...".
\emph{The AAS as a source for control.} In approaches based on capabilities, skills and services, the
AAS describes offered capabilities with their constraints, and skills with their OPC~UA endpoints. A
product's features are matched against the offered capabilities and turned into a manufacturing
sequence~\cite{blumhofer2025enhancing}. How a skill is composed from atomic skills stays inside the
machine's implementation and is not described in the AAS. Stolze et al. make the AAS proactive by
running the bidding of the I4.0 language as BPMN workflows on a process engine, so that a product
selects the resources that offer its required capabilities~\cite{stolze2024proactive}. Albus et al. go
from a product's CAD model to
a bill of process, match it against the resources' capabilities, rebalance the line, and let agents
orchestrate the resources' skills over OPC~UA~\cite{albus2025automated}; they start from the
observation that new processes otherwise need manual integration into the control system. Our earlier
work orchestrated modules from their AASs with behaviour trees~\cite{jensen2026service}. In all of
these, the AAS or the product's model selects, parameterises or calls behaviour that a resource
already has. It does not describe behaviour to be built in the resource's control program.

% Full text: huang2026plug. Abstract only: rooker2007zero.
\emph{Changing IEC~61499 programs online.} The management commands of IEC~61499 create, connect,
parameterise and delete function blocks in a running device~\cite{iec61499}. Reconfiguration services
built on them change part of an application while the rest keeps running~\cite{rooker2007zero}.
Closest to this paper, Huang et al. map the Module Type Package (MTP) onto IEC~61499: a service becomes
a function block with the MTP's service state machine, and an MTP file sent to the runtime is compared
with what runs and turned into management commands~\cite{huang2026plug}. A changed orchestration of a
packaging demonstrator was applied in at most 96\,ms, against 504\,ms with management commands sent one
by one. The orchestration itself is still redrawn by hand. The MTP describes the devices' services, not
what a product needs, and the program is not checked against the file once the change is applied.

% As the outline: nabizada2026capability, schwoerer2026fond. Full text: terzimehic2021service,
% albus2025automated.
\emph{Planning from capabilities.} PDDL is generated from the capability models in the
AAS~\cite{nabizada2026capability}, and policies planned from distributed AAS descriptions run as
behaviour trees~\cite{schwoerer2026fond}. Recipe-driven synthesis decides which controller runs which
service~\cite{terzimehic2021service}, and product-driven line balancing which station does which
step~\cite{albus2025automated}. These plans sequence or assign skills that the resources already
have, for an orchestrator. None turns a plan into a new skill of the resource itself.

\emph{Gap.} No approach makes the AAS the desired state of a module's IEC~61499 program. In such an
approach, a new skill, handcrafted or planned from the module's capabilities, would be checked against
its primitives' conditions and the product's requirement, created in the running program online,
verified by reading the program back, and recorded in the AAS.

% Marks, from the full texts:
% - dorofeev2021orchestrator: a generated basic FB (61499); no online change, check or read-back.
% - terzimehic2021service: a generated IEC 61499 application; SMT satisfiability only (check, in part);
%   how it reaches running PLCs is not described.
% - lindorfer2019user: an external workflow engine; the meta-model rules refuse an invalid model.
% - huang2026plug: management commands online; the runtime is queried and compared with the MTP file
%   before the commands are sent (read back, in part); behaviour contracts are future work.
% - blumhofer2025enhancing: capability constraints checked in production planning (check, in part).
% - stolze2024proactive: BPMN on a process engine; no control program.
% - albus2025automated: validated by simulation (check, in part); no AAS; execution by agents.
\begin{table}[t]
  \centering
  \caption{The closest works: where new behaviour comes from, whether the AAS describes it, whether it
    runs in the controller's IEC~61499 program, whether it is changed online, whether it is checked
    before deployment, and whether the running program is read back and compared with its
    description. \checkmark~yes, (\checkmark)~in part, --~no.}
  \label{tab:related}
  \footnotesize
  \setlength{\tabcolsep}{2.5pt}
  \begin{tabular}{@{}llccccc@{}}
    \hline
    Work & New behaviour from & AAS & 61499 & Online & Check & Read back \\
    \hline
    \cite{dorofeev2021orchestrator} & a given sequence & -- & \checkmark & -- & -- & -- \\
    \cite{terzimehic2021service} & recipe, SMT & -- & \checkmark & -- & (\checkmark) & -- \\
    \cite{lindorfer2019user} & BPMN by the user & -- & -- & -- & \checkmark & -- \\
    \cite{huang2026plug} & MTP file & -- & \checkmark & \checkmark & -- & (\checkmark) \\
    \cite{blumhofer2025enhancing} & product features & \checkmark & -- & -- & (\checkmark) & -- \\
    \cite{stolze2024proactive} & bidding & \checkmark & -- & -- & -- & -- \\
    \cite{albus2025automated} & CAD model & -- & -- & -- & (\checkmark) & -- \\
    Ours & AAS, written or planned & \checkmark & \checkmark & \checkmark & \checkmark & \checkmark \\
    \hline
  \end{tabular}
\end{table}
